\newcommand{\cmark}{\ding{51}}
\newcommand{\xmark}{\ding{55}}
\newcommand{\omark}{\ding{108}}
\chapter{Requirements}
The following shows a detailed overview about which requirements declared in the Requirement Analysis have been implemented.

\section{Must Criteria}
All Must criteria have been fully implemented within the CodeX system.

\begin{longtable}{p{1.2cm} p{11cm} p{1.5cm}}
\caption{Implementation Status of Must Criteria}
\label{tab:must-criteria} \\

\toprule
\textbf{ID} & \textbf{Requirement} & \textbf{Fulfilled} \\
\midrule
\endfirsthead

\toprule
\textbf{ID} & \textbf{Requirement} & \textbf{Fulfilled} \\
\midrule
\endhead

RM1 & The system must provide a web interface that allows natural language to be converted into a snippet of code. The target DSL can be selected from available options and the generated code can be exported as a file. & \cmark \\

RM2 & The system must provide SPICE as one option of its selectable DSLs. & \cmark \\

RM3 & SPICE conversations must include a visualization of the respective circuit. & \cmark \\

RM4 & Every conversation must have a name. & \cmark \\

RM5 & The system must support LLM requests using a Retrieval-Augmented Generation (RAG) approach. Retrieved data must be stored in a secure database that respects intellectual property constraints. & \cmark \\

RM6 & The system must integrate a validation back-end to validate the generated circuit for SPICE conversations. & \cmark \\

RM7 & The system must allow adding and removing separate LLM models together with their OpenAI API URL. When a request is made, the user must be able to select the model with which the request will be performed. & \cmark \\

\bottomrule
\end{longtable}

\section{Should Criteria}
Most Should Criteria were implemented.
RS2 was skipped due to time-pressure.
RS4 was implemented only as a client-side check.
RS8 was partially implemented, as no error message is displayed in that case.

\begin{longtable}{p{1.2cm} p{11cm} p{1.5cm}}
\caption{Implementation Status of Should Criteria}
\label{tab:should-criteria} \\

\toprule
\textbf{ID} & \textbf{Requirement} & \textbf{Fulfilled} \\
\midrule
\endfirsthead

\toprule
\textbf{ID} & \textbf{Requirement} & \textbf{Fulfilled} \\
\midrule
\endhead

RS1 & The circuit visualization for SPICE should be interactive. & \cmark \\

RS2 & The interface of the system should allow users to cancel a pending request. & \xmark \\

RS3 & Users should be able to sign up, log in, and log out using a username (or email) and password. & \cmark \\

RS4 & A strict password complexity policy should be implemented. & \omark \\

RS5 & The system should allow users to securely update their profile data. & \cmark \\

RS6 & The system should provide a secure mechanism for users to recover their accounts in case of lost credentials. & \cmark \\

RS7 & Conversations should be stored with respect to the account that created them and be visible only to that account. & \cmark \\

RS8 & If the LLM fails to generate working code after multiple retries, the system should stop further attempts and display an appropriate error message indicating the failure. & \cmark \\

RS9 & The system should provide a list of all previous conversations and allow users to view them. & \cmark \\

RS10 & Every conversation should be able to be marked as deleted and should not be shown in the overview afterwards. & \cmark \\

\bottomrule
\end{longtable}
\section{Could Criteria}

Most Could Criteria were implemented.
RC1 was skipped due to time-pressure.
RC7 was implemented only as a client-side check.
RC12 was partially implemented, as no error message is displayed in that case.

\begin{longtable}{p{1.2cm} p{11cm} p{1.5cm}}
\caption{Implementation Status of Could Criteria}
\label{tab:could-criteria} \\

\toprule
\textbf{ID} & \textbf{Requirement} & \textbf{Fulfilled} \\
\midrule
\endfirsthead

\toprule
\textbf{ID} & \textbf{Requirement} & \textbf{Fulfilled} \\
\midrule
\endhead

RC1 & The system could integrate a simulation back-end to simulate the generated circuit for SPICE conversations. & \cmark \\

RC2 & The system could accept follow-up instructions on conversations. & \cmark \\

RC3 & The system could provide more DSL options than only SPICE. & \cmark \\

RC4 & Every account could be associated with the role Admin or Member. Each role could have a configurable upper bound of requests per month set as a server parameter. & \cmark \\

RC5 & The role request upper bounds could be configured by administrators via the web interface. & \cmark \\

RC6 & The web interface could be offered in multiple languages, at least German and English. & \cmark \\

RC7 & If supported by the DSL, the generated code output could be prettified using syntax highlighting. & \xmark \\

RC8 & The front-end could allow users to pin a conversation so that it remains at the top of the conversation list. & \cmark \\

RC9 & During registration, the user’s email could be confirmed via a verification link or code. Registration would not be completed until confirmation. Changing the email would automatically trigger a new confirmation. & \cmark \\

RC10 & The system could support OAuth 2.0 authentication to allow login via trusted providers such as Google. & \cmark \\

RC11 & The software could provide a dark mode as an alternative theme. & \cmark \\

RC12 & The system could integrate with cloud storage providers to save or share generated code snippets. & \xmark \\

RC13 & The application could allow users to search conversations by name. & \cmark \\

RC14 & The front-end could indicate to the user that a request is currently being processed. & \cmark \\

RC15 & The software could allow users to download generated code as a file. & \cmark \\

RC16 & The software could provide admin-only analytics on usage (e.g., daily active users) not visible to regular users. & \xmark \\

RC17 & Users could upload and display a profile picture associated with their account. & \cmark \\

\bottomrule
\end{longtable}

\section{QUality Criteria}
More than half of the quality criteria is inherently met.
Criteria Q30, Q90 and Q130 were found to be out-of-scope of the implementation and are the responsibility of the system administrator.
Q50 and Q110 were not implemented due to time pressure.
Q80 was found to not be implementable due to 3rd-party libraries producing non-preventable errors even during normal operation.

\begin{longtable}{p{1.2cm} p{11cm} p{1.5cm}}
\caption{Implementation Status of Could Criteria}
\label{tab:could-criteria} \\

\toprule
\textbf{ID} & \textbf{Requirement} & \textbf{Fulfilled} \\
\midrule
\endfirsthead

\toprule
\textbf{ID} & \textbf{Requirement} & \textbf{Fulfilled} \\
\midrule
\endhead

Q10 & Database injection must not be possible. & \cmark \\

Q20 & No personally identifiable information (PII) must be shared with external services. & \cmark \\

Q30 & The code generation must not take longer than 30 seconds. In case it takes longer, that
generation is counted as failed. & \omark \\

Q40 & All user passwords should only be persisted using a secure, one-way encryption algo-
rithm. & \cmark \\

Q50 & The system should limit concurrent users to prevent overloading the server. & \xmark \\

Q60 & All server-side logging should be performed using a logging framework to ensure struc-
tured logs, maintainability and effective debugging. & \cmark \\

Q70 & The front-end should handle all operations without runtime errors. & \cmark \\

Q80 & The front-end should not produce any errors in the browser console during normal
operation. & \xmark \\

Q90 & All communication between front-end and back-end should be encrypted. & \omark \\

Q100 & The system could be deployable using containerization. & \cmark \\

Q110 & The system could allow integration with CI/CD pipelines to automatically test or
deploy generated code. & \xmark \\

Q120 & The interface of the system should be responsive on the display sizes Mobile Phone,
Tablet and Desktop. It should also be accessible according to the Material Accessibility Guidelines. & \cmark \\

Q130 & The LLM and Database should be stored on a secure server and can only be accessed
by authorized actors. & \omark \\

\bottomrule
\end{longtable}
